\newpage

\addcontentsline{toc}{section}{Введение}
\section*{\centering{\uppercase{Введение}}}

Фрактальная графика относится к области визуализации сложных математических моделей и активно используется в цифровом искусстве, дизайне и научной визуализации. Одним из наиболее выразительных направлений является генерация фрактального «пламени», предложенная С. Дрейвсом, где изображение формируется итеративными нелинейными преобразованиями и последующей тональной обработкой плотности \cite{fractalflame}. Благодаря богатству форм и высокой детализации такие изображения востребованы как в творческих проектах, так и в обучении математическому моделированию.

Современные инструменты генерации фракталов чаще представлены настольными приложениями и ориентированы на локальную работу пользователя. При этом растёт спрос на веб-сервисы, обеспечивающие доступ с разных устройств, хранение истории и возможность публикации результатов. Для решения этих задач требуется система, сочетающая удобный пользовательский интерфейс, вычислительно ёмкий рендеринг и надёжное хранение артефактов.

Целью курсового проекта является создание распределённой системы генерации фракталов с веб-интерфейсом, поддержкой аутентификации пользователей и хранением истории созданных работ. Архитектура системы построена на принципах микросервисов, что облегчает масштабирование и независимое развитие компонентов \cite{microservices}. Разработка выполнена в соответствии с принципами объектно-ориентированного проектирования и ориентирована на практическое применение современных инструментов сборки и развёртывания.

Для достижения поставленной цели были решены следующие задачи:
\begin{itemize}
  \item анализ существующих решений и выявление их сильных и слабых сторон;
  \item формирование функциональных и нефункциональных требований к системе;
  \item проектирование микросервисной архитектуры и модели данных;
  \item выбор технологий реализации, сборки и развёртывания;
  \item разработка и тестирование ключевых сценариев работы пользователей.
\end{itemize}

\newpage

\section{Анализ аналогов}

Для определения требований и позиционирования проекта были рассмотрены наиболее известные решения в области генерации фрактального пламени. Сравнение выполнено по критериям: доступность (desktop/web), качество рендеринга, удобство настройки параметров, наличие истории и механизмов публикации результатов.

\subsection{Apophysis}

\begin{figure}[H]
  \centering
  \includegraphics[width=\textwidth]{apophysis.png}
  \caption{Скриншот пользовательского интерфейса Apophysis с несколькими открытыми окнами}
\end{figure}

Apophysis --- один из наиболее известных редакторов фрактального пламени. Программа распространяется как open-source и предоставляет богатый набор инструментов для ручной настройки трансформаций, палитр и параметров рендера \cite{apophysis}. К достоинствам можно отнести простоту использования и обилие обучающих материалов. Основные ограничения связаны с настольным форматом, отсутствием встроенного хранения истории и невозможностью работы в распределённой среде.

Дополнительно стоит отметить, что основная работа в Apophysis строится вокруг ручного подбора вариаций и параметров, поэтому пользователь часто вынужден сохранять множество промежуточных файлов. Отсутствие центрального хранилища и единой учётной записи делает обмен результатами и совместную работу затруднительными.

\subsection{JWildfire}

\begin{figure}[H]
  \centering
  \includegraphics[width=\textwidth]{jwildfire.png}
  \caption{Скриншот пользовательского интерфейса JWildfire}
\end{figure}

JWildfire --- мощный кроссплатформенный редактор фрактального пламени с большим набором вариаций и встроенными генераторами случайных фракталов \cite{jwildfire}. Решение предоставляет широкие возможности, но требует значительного времени на освоение и ориентировано на опытных пользователей. Встроенные механизмы совместной работы и веб-публикации отсутствуют, а хранение результатов полностью локально.

Сильной стороной JWildfire является гибкость и глубина настроек, что делает его удобным для экспериментов и профессионального использования. Однако сложный интерфейс и необходимость локальной установки снижают доступность для начинающих пользователей и не решают задачу быстрого доступа к проектам с разных устройств.

\subsection{Chaotica}

\begin{figure}[H]
  \centering
  \includegraphics[width=\textwidth]{chaotica.png}
  \caption{Скриншот анимации фрактала, создаваемой в Chaotica}
\end{figure}

Chaotica --- коммерческое приложение, ориентированное на получение высококачественных изображений и ускоренный рендеринг (включая GPU) \cite{chaotica}. Преимущества: высокая скорость и качество, развитый инструментарий анимации. Недостатки: закрытость, платная модель распространения и ориентация на локальную работу без интегрированного облачного хранения и пользовательских профилей.

В практической эксплуатации продукт удобен для студийных задач и финального качества рендера, однако ориентирован на автономное использование. Это ограничивает возможность модификаций, интеграции с веб-архивом и построения персонализированного профиля пользователя.

\subsection{Сравнительный анализ}

\begin{table}[H]
  \centering
  \caption{Сравнение аналогов по ключевым критериям}
  \begin{tabular}{|p{3.6cm}|p{5.0cm}|p{6.0cm}|}
    \hline
    Продукт   & Сильные стороны                                                          & Ограничения                                                             \\
    \hline
    Apophysis & Простота использования, open-source, богатая база примеров               & Только desktop, нет истории и веб-доступа                               \\
    \hline
    JWildfire & Большое количество вариаций, кросс-платформенность, мощный набор функций & Высокая сложность интерфейса, локальное хранение, отсутствие веб-режима \\
    \hline
    Chaotica  & Высокое качество рендера, поддержка GPU, развитая анимация               & Платная лицензия, закрытость, отсутствие облачной истории               \\
    \hline
  \end{tabular}
\end{table}

Анализ показал, что существующие решения преимущественно ориентированы на локальный рабочий процесс и не обеспечивают удобной работы с историей и публикацией результатов. Это определило требования к проектируемой системе, в частности необходимость веб-интерфейса, учётных записей и долговременного хранения рендеров.

Следовательно, ключевым конкурентным преимуществом будущей системы становится доступность из браузера и единое пространство проектов, где результаты сохраняются и доступны для повторного использования.

\section{Формирование функциональных требований}

На основе анализа аналогов сформированы функциональные требования к системе генерации фракталов:
\begin{itemize}
  \item регистрация, вход и обновление пользовательской сессии с использованием JWT-токенов;
  \item создание задачи генерации фрактала с передачей размеров изображения, количества итераций, уровня симметрии, гамма-коррекции, набора функций и аффинных преобразований;
  \item получение списка доступных функций фрактального пламени;
  \item отображение прогресса генерации и промежуточного предпросмотра;
  \item получение списка ранее созданных фракталов с пагинацией, сортировкой и фильтрацией по пользователю;
  \item просмотр отдельного фрактала, его параметров и итогового изображения;
  \item отображение главной страницы с формой генерации и feed последних сгенерированных фракталов с подписями пользователей;
  \item отображение страницы профиля с редактированием пользовательской информации и историей созданных фракталов;
  \item сохранение параметров генерации в PostgreSQL и сохранение PNG-изображений в S3-совместимом хранилище MinIO;
  \item редактирование базовых данных профиля пользователя.
\end{itemize}

Сформированы также нефункциональные требования:
\begin{itemize}
  \item производительность: приемлемое время рендера при стандартных параметрах качества;
  \item масштабируемость: возможность независимо увеличивать количество экземпляров gateway и generator без изменения клиентского приложения;
  \item надёжность: сохранность истории генераций в PostgreSQL и изображений в MinIO при перезапуске контейнеров;
  \item безопасность: stateless-аутентификация через JWT, передача токена в HTTP-заголовке и проксирование его во внутренние gRPC-вызовы;
  \item переносимость: развёртывание всех компонентов через Docker Compose;
  \item сопровождаемость: разделение frontend, API gateway, IAM и generator по самостоятельным зонам ответственности;
  \item ограничение вычислительной нагрузки: проверка максимальной ширины, высоты и числа итераций на стороне generator.
\end{itemize}

\section{Анализ литературы}

\subsection{Принципы объектно-ориентированного проектирования}

Объектно-ориентированный подход опирается на понятия инкапсуляции, наследования, полиморфизма и абстракции. Эти принципы позволяют моделировать предметную область в виде взаимосвязанных классов, обеспечивая повторное использование и расширяемость кода \cite{ooaod_booch}. Для разрабатываемой системы это особенно важно, поскольку логика генерации фракталов включает множество разновидностей преобразований, параметров рендера и вариаций, которые удобно представлять в виде иерархий объектов.

В проекте объектные модели описывают параметры рендера, состояние проекта и пользовательский контекст, что позволяет разделить вычислительную и интерфейсную части. Это упрощает развитие функциональности и делает код более устойчивым к изменениям.

\subsection{Принципы SOLID и качество кода}

Принципы SOLID ориентированы на снижение связности и повышение устойчивости к изменениям. В частности, разделение ответственности и открытость к расширению позволяют добавлять новые алгоритмы генерации и модули без переработки существующего кода \cite{solid_martin}. Эти принципы были использованы при проектировании сервисов, чтобы обеспечивать независимость компонентов и удобство тестирования.

На практике это выразилось в выделении отдельных интерфейсов для рендера, хранения и пользовательских операций, а также в уменьшении числа прямых зависимостей между слоями. Такой подход упростил модульное тестирование и локализацию ошибок.

\subsection{Шаблоны проектирования}

В литературе по объектно-ориентированному проектированию выделяются шаблоны, упрощающие решение типовых инженерных задач \cite{design_patterns}. Для проекта генерации фрактального пламени особенно важны шаблоны, которые позволяют отделить изменяемые алгоритмы от стабильной инфраструктуры приложения.

Шаблон Strategy применим к вариационным функциям фрактального пламени. Каждая функция преобразует точку по собственному правилу, но внешний код генерации может работать с общим интерфейсом функции. Благодаря этому добавление новой вариации не требует переписывать основной цикл рендера. Достаточно реализовать новую модель функции и зарегистрировать её в фабрике или построителе функций.

Шаблон Factory Method полезен при создании доменных объектов на основе входных DTO или protobuf-сообщений. Внешний контракт передаёт имя функции и числовые параметры, а внутри generator должен появиться объект конкретного класса. Вынесение этой логики из сервиса генерации снижает связность и делает ошибочные имена функций отдельным контролируемым сценарием.

Идея Builder применима к созданию задачи генерации. Задача включает размеры изображения, число итераций, симметрию, гамму, набор функций, набор аффинных преобразований и пользовательский контекст. Последовательная сборка такой конфигурации удобнее, чем передача большого количества несвязанных параметров между методами.

\subsection{Моделирование предметной области}

Для описания предметной области генерации фракталов применяются подходы доменно-ориентированного проектирования. Они позволяют выделить ключевые сущности и описать их жизненный цикл \cite{ddd_evans}. В данном проекте основными сущностями являются пользователь, фрактал, вариационная функция, аффинное преобразование, задача генерации и изображение результата.

Сущность фрактала хранит параметры, достаточные для повторного отображения результата: размеры, количество итераций, уровень симметрии, гамму, список функций и список аффинных преобразований. Изображение при этом не смешивается с параметрами: в реляционной базе хранится только имя объекта, а бинарное содержимое помещается в объектное хранилище. Такое разделение поддерживает принцип единственной ответственности на уровне данных.

Пользовательская сущность относится к другому bounded context --- IAM. Она не должна зависеть от деталей генерации фракталов. Поэтому generator хранит только внешний идентификатор пользователя, полученный из JWT, а полная информация о профиле остаётся в IAM. Это уменьшает связность микросервисов и упрощает изменение профиля без миграций в сервисе генерации.

\subsection{Слоистая организация приложения}

Серверные сервисы проекта разделяются на несколько слоёв: транспортный, сервисный, доменный и инфраструктурный. Транспортный слой принимает HTTP или gRPC-запросы и отвечает за преобразование данных из внешнего формата. Сервисный слой координирует сценарии использования: запуск генерации, получение истории, чтение изображения, обновление пользователя. Доменный слой содержит модели фрактала, функций и параметров. Инфраструктурный слой работает с PostgreSQL, MinIO и внешними библиотеками.

Такое разделение важно для тестируемости. Например, алгоритм генерации можно проверять отдельно от gRPC-сервера, а мапперы DTO можно тестировать без запуска базы данных. Кроме того, слой gateway не обязан знать внутреннее устройство generator: он опирается на контракт и выполняет роль адаптера между HTTP-клиентом и внутренними сервисами.

\subsection{Контракты и переносимость моделей}

В распределённой системе объектная модель существует сразу в нескольких формах: как доменная модель внутри сервиса, как схема базы данных, как protobuf-сообщение для gRPC и как JSON-модель во внешнем HTTP API. Если эти формы изменяются вручную и независимо, быстро появляются расхождения между клиентом и сервером.

Для снижения этого риска используются формальные контракты. gRPC-контракт generator описывает внутренние сообщения и сервисные методы, а OpenAPI-спецификация описывает внешний HTTP API gateway. Кодогенерация по этим контрактам позволяет получать типы, интерфейсы и модели из единого источника. Это не отменяет ручное проектирование доменной модели, но делает границы между сервисами более проверяемыми.

\section{Проектирование системы}

\subsection{Общая архитектура}

\begin{figure}[H]
  \centering
  \includegraphics[width=\textwidth]{architecture.png}
  \caption{Диаграмма архитектуры}
\end{figure}

Проект реализуется как распределённая система из четырёх компонентов:
\begin{itemize}
  \item frontend --- клиентское приложение на TypeScript, Vue, Tailwind CSS и Axios;
  \item API gateway --- единая HTTP-точка входа для клиентского приложения, реализованная на Java и Spring;
  \item IAM --- сервис управления пользователями и аутентификацией, реализованный на Go;
  \item generator --- сервис генерации и хранения фракталов, реализованный на Java и Spring.
\end{itemize}

Frontend взаимодействует только с API gateway по HTTP. Gateway скрывает внутреннюю структуру backend-части и обращается к IAM и generator по gRPC \cite{grpc}. Такое разделение позволяет независимо развивать пользовательский интерфейс, аутентификацию и вычислительную часть. Внешний контракт остаётся ориентированным на браузерное приложение, а внутренние сервисы могут использовать более строгие protobuf-контракты.

Зоны ответственности компонентов представлены в таблице.

\begin{table}[H]
  \centering
  \caption{Ответственность компонентов системы}
  \begin{tabularx}{\textwidth}{|p{3.3cm}|X|X|}
    \hline
    Компонент   & Назначение                                                            & Взаимодействие                                                                    \\
    \hline
    Frontend    & Отображение интерфейса, ввод параметров, просмотр галереи и профиля   & HTTP-запросы к gateway, получение событий прогресса и изображений                 \\
    \hline
    API gateway & Единая точка входа и адаптер между HTTP API и внутренними сервисами   & HTTP с frontend, gRPC-вызовы IAM и generator, передача JWT во внутренние metadata \\
    \hline
    IAM         & Регистрация, вход, обновление токенов и хранение профиля пользователя & gRPC API для gateway, хранение пользовательской записи в PostgreSQL               \\
    \hline
    Generator   & Генерация фрактального пламени, история работ, хранение изображений   & gRPC API для gateway, PostgreSQL для параметров, MinIO для PNG-файлов             \\
    \hline
  \end{tabularx}
\end{table}

\subsection{Frontend}

Frontend отвечает за пользовательский сценарий работы с системой. Он реализуется как одностраничное приложение, в котором выделены две основные страницы: главная страница генерации и страница профиля. Главная страница объединяет форму параметров фрактала, блок текущего состояния задачи и feed последних сгенерированных работ. В feed отображается миниатюра результата, краткие параметры и подпись пользователя, создавшего фрактал. Галерея примеров результатов должна быть приведена в приложении~А.

Страница профиля решает другую группу задач: редактирование пользовательской информации и просмотр личной истории генераций. В верхней части страницы размещается форма изменения имени и изображения профиля, а ниже --- список фракталов текущего пользователя. Такое разделение интерфейса поддерживает два режима работы: быстрый творческий цикл на главной странице и управление личным архивом на странице профиля.

Клиентский код не обращается напрямую к IAM или generator. Все запросы идут через API gateway, что упрощает хранение базового URL, обработку ошибок и передачу JWT. HTTP-контракт описывается OpenAPI-спецификацией \cite{openapi}; на стороне frontend она используется как источник типизированных моделей и клиентских функций. Благодаря этому изменение структуры запроса или ответа быстрее обнаруживается на этапе сборки.

% TODO: вставить скриншот главной страницы сайта: форма генерации фрактала, блок прогресса/результата и feed последних сгенерированных фракталов с подписями пользователей.
% TODO: вставить скриншот страницы профиля: форма изменения информации пользователя и история сгенерированных фракталов.

\subsection{API gateway}

API gateway является границей backend-части. Он предоставляет frontend HTTP-интерфейс в формате JSON и скрывает внутренние gRPC-сервисы. Для работы с фракталами предусмотрены маршруты получения списка фракталов, получения фрактала по идентификатору, запуска генерации, подписки на события задачи, получения списка функций и получения PNG-изображения. Для пользовательских сценариев предусмотрены маршруты регистрации, входа, обновления токена, получения и изменения профиля.

В gateway внешний API описан в файле \texttt{openapi.yaml}. По этой спецификации OpenAPI Generator формирует Java-модели и API-интерфейсы, которые затем реализуются контроллерами. Это снижает вероятность расхождения между документацией, серверной реализацией и клиентским кодом.

Внутри gateway слой HTTP-контроллеров отделён от слоя вызовов downstream-сервисов. Контроллер \texttt{FractalsController} принимает запросы на маршрутах \texttt{/api/v1/fractals}, \texttt{/api/v1/fractals/gen}, \texttt{/api/v1/fractals/functions} и \texttt{/api/v1/fractals/images/\{name\}}. Сервисный слой использует gRPC-stub'ы generator: блокирующий stub применяется для коротких операций, а асинхронный stub --- для подписки на состояние задачи.

Gateway также отвечает за передачу пользовательского контекста. JWT приходит от frontend в HTTP-заголовке \texttt{Authorization}, после чего gateway передаёт токен во внутренние gRPC metadata. В результате generator может связать созданный фрактал с пользователем, не имея прямого HTTP-взаимодействия с браузером.

\subsection{IAM}

IAM выделен в отдельный микросервис, чтобы логика пользователей и токенов не смешивалась с генерацией изображений. Сервис предоставляет внутренний gRPC API для регистрации, входа, обновления токена, получения пользователя и изменения профиля. Внешние HTTP-маршруты остаются на gateway, поэтому frontend не зависит от конкретной реализации IAM.

Аутентификация построена на JWT \cite{jwt_rfc}. После успешного входа IAM выдаёт access token и refresh token. Access token используется frontend при обращении к gateway, а refresh token позволяет обновлять пользовательскую сессию без повторного ввода пароля. Токен содержит идентификатор пользователя, который downstream-сервисы используют как связку между действиями пользователя и данными системы.

Для IAM предусмотрена отдельная модель пользователей. В рамках текущих требований она состоит из одной таблицы \texttt{users}, где хранятся идентификатор пользователя, имя пользователя, хеш пароля и ссылка на изображение профиля или имя объекта в хранилище. Идентификатор пользователя используется как subject JWT и передаётся downstream-сервисам для связывания пользовательских действий с историей генераций.

% TODO: вставить ER-диаграмму IAM: users(id, username, password_hash, image), уникальное ограничение на username.

\subsection{Generator}

Generator отвечает за вычислительно тяжёлую часть системы: создание фрактального пламени, хранение параметров генерации, хранение результата и выдачу истории работ. Внешние клиенты не обращаются к нему напрямую; доступ выполняется через API gateway по внутреннему gRPC-контракту.

Основной контракт generator описан в файле \texttt{generator.proto}. Сервис \texttt{Fractals} предоставляет следующие операции:
\begin{itemize}
  \item \texttt{GetAll} --- получение списка фракталов с пагинацией, сортировкой и фильтрацией по \texttt{user\_id};
  \item \texttt{GetById} --- получение одного фрактала по идентификатору;
  \item \texttt{Generation} --- создание задачи генерации;
  \item \texttt{SubscribeToTask} --- серверный поток состояния задачи;
  \item \texttt{GetFunctions} --- получение доступных имён вариационных функций;
  \item \texttt{GetImage} --- чтение изображения из S3-совместимого хранилища.
\end{itemize}

Запрос генерации содержит ширину и высоту изображения, число итераций, уровень симметрии, параметр гамма-коррекции, список вариационных функций и список аффинных преобразований. Перед постановкой задачи в очередь generator проверяет ограничения параметров: максимальную ширину, высоту, число итераций, положительный уровень симметрии, наличие хотя бы одной функции и одного набора аффинных параметров.

Сервис также содержит каталог поддерживаемых функций фрактального пламени. По коду generator доступны функции \texttt{linear}, \texttt{sinusoidal}, \texttt{spherical}, \texttt{swirl}, \texttt{horseshoe}, \texttt{polar}, \texttt{disc}, \texttt{heart} и \texttt{eyefish}. Каждая функция представлена отдельной моделью, что упрощает добавление новых вариаций без изменения внешнего API.

\begin{figure}[H]
  \centering
  \includegraphics[width=\textwidth]{generator-dbml.png}
  \caption{ER-диаграмма сервиса генерации}
\end{figure}

Модель данных generator строится вокруг истории созданных фракталов. По миграциям Liquibase используются три таблицы: \texttt{fractals}, \texttt{functions} и \texttt{affine\_params}. Таблица \texttt{fractals} хранит общие параметры результата: идентификатор, идентификатор пользователя, дату создания, размеры изображения, число итераций, уровень симметрии, гамму и имя файла изображения. Таблица \texttt{functions} хранит набор вариационных функций и веса для конкретного фрактала. Таблица \texttt{affine\_params} хранит цвет и коэффициенты аффинного преобразования \texttt{a}, \texttt{b}, \texttt{c}, \texttt{d}, \texttt{e}, \texttt{f}.

Связи имеют тип «один ко многим»: один фрактал содержит несколько функций и несколько наборов аффинных параметров. Для дочерних таблиц настроено каскадное удаление по внешнему ключу \texttt{fractal\_id}, поэтому удаление фрактала не оставляет несвязанных параметров. Для ускорения выборок созданы индексы по \texttt{fractals.user\_id}, \texttt{functions.fractal\_id} и \texttt{affine\_params.fractal\_id}.

Метаданные фракталов хранятся в PostgreSQL, а бинарные PNG-файлы вынесены в MinIO --- локальное S3-совместимое объектное хранилище \cite{minio}. Такое разделение уменьшает нагрузку на реляционную базу данных и упрощает выдачу файлов. Generator использует AWS S3 SDK, создаёт bucket \texttt{images} при запуске, сохраняет PNG с типом \texttt{image/png} и возвращает имя файла, которое затем хранится в таблице \texttt{fractals}.

При запросе изображения gateway вызывает gRPC-метод \texttt{GetImage}; generator читает объект из MinIO и возвращает байтовое содержимое. Для frontend это остаётся обычным HTTP-ответом с изображением, а внутренняя реализация хранения скрыта за gateway.

\subsection{Процесс генерации фрактального пламени}

Генерация фрактального пламени строится как итерационный процесс построения изображения по набору аффинных преобразований и вариационных функций. Пользователь передаёт размеры изображения, количество итераций, уровень симметрии, гамму, веса функций и коэффициенты аффинных преобразований. Эти параметры сохраняются как часть доменной модели фрактала, после чего на их основе создаётся объект задачи \texttt{GenerationTask}. Фрагменты кода генератора и сервиса генерации должны быть приведены в приложении~Б.

Внутри рендера выбирается начальная случайная точка и случайный цвет. На каждой итерации выбирается одно аффинное преобразование и одна вариационная функция. Аффинное преобразование переносит точку в новое положение, а вариационная функция применяет нелинейное преобразование. Затем точка дополнительно обрабатывается с учётом уровня симметрии: для каждого симметричного поворота рассчитываются координаты пикселя на плоскости изображения.

Первые итерации используются как прогрев и не записываются в изображение, чтобы начальное случайное состояние меньше влияло на результат. После прогрева координаты точки преобразуются в координаты пикселя, цвет смешивается с цветом текущего аффинного преобразования, а счётчик посещений пикселя увеличивается. Чем чаще точка попадает в область изображения, тем сильнее эта область проявляется после нормализации.

Для ускорения вычислений generator делит работу между несколькими потоками. Класс \texttt{Generator} создаёт пул потоков, а каждый экземпляр \texttt{Render} выполняет свою часть итераций над общей моделью изображения. Доступ к пиксельной матрице синхронизируется, поскольку несколько потоков могут одновременно обновлять один и тот же результат. После завершения потоков выполняется гамма-коррекция: счётчики посещений переводятся в логарифмическую нормализацию, после чего цвет пикселей усиливается с учётом параметра \texttt{gamma}.

Результат генерации сохраняется в MinIO как PNG-файл, а в PostgreSQL остаётся ссылка на объект и параметры, достаточные для отображения истории. Такой подход делает результат воспроизводимым на уровне параметров и одновременно не перегружает реляционную базу бинарными данными.

\subsection{Передача прогресса генерации}

Генерация фрактала может занимать заметное время, поэтому пользователь не должен ждать завершения HTTP-запроса без обратной связи. Для этого генерация реализована как асинхронная задача. Метод \texttt{Generation} в generator сохраняет параметры, создаёт \texttt{GenerationTask}, добавляет её в очередь и сразу возвращает начальное состояние задачи с идентификатором фрактала и текущим прогрессом.

Внутри \texttt{GeneratationService} используется очередь \texttt{BlockingQueue} и набор ожидающих задач \texttt{ConcurrentHashMap}. При запуске сервиса создаётся фиксированный пул worker-потоков. Каждый worker берёт задачу из очереди, вызывает \texttt{generator.process(task)}, после завершения сохраняет PNG-изображение в MinIO, обновляет запись фрактала в PostgreSQL и удаляет задачу из списка ожидающих.

Прогресс считается в \texttt{GenerationTask}. Внутри задачи хранится счётчик выполненных итераций \texttt{AtomicLong}; поток рендера периодически увеличивает его пакетами итераций. Метод \texttt{toTaskState()} формирует protobuf-сообщение \texttt{TaskState}: идентификатор фрактала, число прогресса от 0 до 1 и base64-предпросмотр текущего изображения.

Рассылка обновлений выполняется через scheduler Spring. Метод \texttt{tasksUpdate()} помечен аннотацией \texttt{@Scheduled(fixedDelayString = "\$\{app.tasks-update-delay\}")}; в конфигурации значение задержки задано как \texttt{1s}. Scheduler проходит по ожидающим задачам и вызывает \texttt{sendUpdate(task)}. Если для задачи есть подписчик, сервис отправляет очередное \texttt{TaskState} в gRPC stream.

Gateway подписывается на gRPC-метод \texttt{SubscribeToTask} асинхронным stub'ом. Каждое сообщение \texttt{TaskState}, пришедшее от generator, преобразуется в модель внешнего API и отправляется браузеру через \texttt{SseEmitter}. Для frontend это выглядит как поток \texttt{text/event-stream}, то есть Server-Sent Events \cite{mdn_sse}. Когда прогресс достигает 1, gateway завершает SSE-поток. Если gRPC stream завершается ошибкой, \texttt{SseEmitter} также завершается с ошибкой, чтобы клиент мог показать пользователю корректное состояние.

% TODO: вставить sequence diagram запуска генерации и передачи прогресса: frontend -> gateway -> generator -> scheduler -> gRPC stream -> SSE.

\subsection{Развёртывание}

Развёртывание выполняется через Docker Compose \cite{docker_compose}. Каждый сервис запускается в отдельном контейнере, а инфраструктурные зависимости поднимаются рядом с приложением. Для generator используются PostgreSQL и MinIO; контейнер generator получает адрес PostgreSQL через переменную \texttt{POSTGRES\_URL}, адрес MinIO через \texttt{MINIO\_URL}, а секрет JWT через \texttt{JWT\_SIGNING\_KEY}. Gateway получает адрес generator через \texttt{GENERATOR\_URL}; аналогично для полной конфигурации задаётся адрес IAM.

Compose-сценарий содержит healthcheck'и для базы данных, MinIO, gateway и generator. Это снижает вероятность запуска приложения до готовности зависимостей и делает локальное окружение воспроизводимым. На сервере запускается тот же Compose-проект, что используется при локальной проверке; отличаются только значения переменных окружения, доменное имя и настройки обратного прокси.

Финальная версия системы разворачивается в Docker-контейнерах на VPS. Для размещения серверной части использовался хостинг-провайдер Jino.ru, выбранный из-за SSH-доступа, возможности управления сервером и интеграции доменов с SSL-сертификатами \cite{jino}. В составе развёртывания применяется обратный прокси для маршрутизации запросов и разграничения доступа к сервисам.

\begin{figure}[H]
  \includegraphics[width=\textwidth]{jino.png}
  \caption{Страница сервера на jino.ru}
\end{figure}

\subsection{Технологии используемые для решения поставленных задач}

\subsubsection{Язык программирования Java}

\begin{figure}[H]
  \centering
  
\includegraphics[width=0.5\textwidth]{java.png}
  \caption{Java}
\end{figure}

Язык программирования Java является объектно-ориентированным языком, разработанным компанией Sun Microsystems и в настоящее время поддерживаемым корпорацией Oracle. Java обеспечивает высокий уровень абстракции, платформенную независимость и безопасность, что делает его одним из самых популярных языков в мире.

Одним из основных преимуществ Java является его независимость от платформы, достигаемая благодаря использованию виртуальной машины Java (JVM), которая исполняет байт-код, одинаковый для всех операционных систем. Это свойство позволяет создавать переносимые приложения без необходимости модификации исходного кода.

Java поддерживает основные принципы объектно-ориентированного программирования, такие как наследование, инкапсуляция, полиморфизм и абстракция. Эти особенности обеспечивают высокую модульность и повторное использование кода. \cite{java}

В рамках проекта Java используется в двух сервисах: API gateway и generator. Для этих компонентов язык выбран благодаря следующим преимуществам:
\begin{itemize}
  \item наличие развитой экосистемы для HTTP API, gRPC, работы с базами данных и контейнеризированного запуска;
  \item развитая экосистема, включающая библиотеки и фреймворки, ускоряющие процесс разработки;
  \item встроенная система управления памятью с автоматическим сборщиком мусора, упрощающая разработку и снижающая вероятность ошибок, связанных с управлением ресурсами;
  \item высокая читаемость и строгая типизация, что способствует написанию надёжного и легко сопровождаемого кода.
\end{itemize}

В gateway Java применяется для обработки HTTP-запросов, настройки фильтра JWT и вызова gRPC-сервисов. В generator язык используется для реализации алгоритма фрактального пламени, очереди задач, сохранения метаданных и работы с S3-совместимым хранилищем.

\subsubsection{Фреймворк Spring}

\begin{figure}[H]
  \centering
  
\includegraphics[width=0.75\textwidth]{spring.png}
  \caption{Spring Framework}
\end{figure}

Spring --- это мощный и гибкий фреймворк с открытым исходным кодом для разработки корпоративных и веб-приложений на языке Java. Благодаря своей модульной структуре, высокой производительности и удобству масштабирования, Spring широко применяется при создании распределённых и высоконагруженных систем, включая веб-сервисы и многопользовательские онлайн-платформы.

Одной из ключевых особенностей Spring является внедрение зависимостей (Dependency Injection), которое способствует слабой связанности компонентов приложения и облегчает их тестирование, повторное использование и расширение. Фреймворк построен по принципу инверсии управления (Inversion of Control), где управление жизненным циклом объектов осуществляется специальным контейнером.

В проекте Spring используется в gateway и generator. Gateway принимает HTTP-запросы от frontend, а generator предоставляет gRPC API для генерации фракталов и доступа к истории.

В процессе разработки были использованы следующие модули Spring:
\begin{itemize}
  \item Spring Boot --- модуль, предназначенный для упрощения конфигурации и быстрой инициализации проекта. Он предоставляет встроенный сервер и автоматическую настройку компонентов, что позволило сосредоточиться на реализации логики приложения;

  \item Spring Web --- модуль, реализующий обработку HTTP-запросов, маршрутизацию и построение REST API gateway;

  \item Spring gRPC --- модуль для реализации gRPC-сервера в generator и gRPC-клиента в gateway;

  \item Spring Data JPA --- модуль для работы с PostgreSQL через объектно-реляционное отображение;

  \item Spring Security --- модуль, обеспечивающий обработку JWT в HTTP-запросах gateway и gRPC metadata generator.
\end{itemize}

Использование Spring позволило добиться высокой скорости обработки запросов, надёжности и масштабируемости серверной части. Благодаря модульной архитектуре фреймворка, разработка велась по принципу разделения ответственности, что улучшило структуру проекта и упростило поддержку и развитие функционала. \cite{spring}

Дополнительно применяются встроенные механизмы валидации данных, централизованная обработка ошибок и healthcheck'и, что обеспечивает единый формат ответов и предсказуемость поведения сервисов.

\subsubsection{Язык Go}

Язык Go используется для реализации IAM-сервиса. Для сервиса аутентификации важны быстрый запуск, простая модель конкурентного выполнения, компактная поставка и удобная работа с сетевыми сервисами. Go хорошо подходит для небольших backend-компонентов, которые предоставляют API, работают с базой данных и не требуют тяжёлой runtime-среды \cite{go}.

В рамках проекта IAM отвечает за регистрацию пользователей, проверку пароля, выпуск JWT и обновление сессии. Выделение этого сервиса на Go позволяет развивать пользовательскую часть независимо от Java-сервисов генерации.

\subsubsection{Frontend: TypeScript, Vue, Tailwind CSS и Axios}

Frontend реализуется как одностраничное веб-приложение. TypeScript добавляет статическую типизацию поверх JavaScript и помогает согласовать модели запросов и ответов с OpenAPI-спецификацией \cite{typescript}. На основе OpenAPI-контракта формируются типы и клиентские функции для обращения к gateway, поэтому изменения HTTP API быстрее обнаруживаются на этапе сборки frontend. Vue выбран как фреймворк для построения реактивного интерфейса: он упрощает разбиение экрана на компоненты, управление состоянием формы генерации и обновление истории фракталов \cite{vue}.

Tailwind CSS применяется для быстрой сборки адаптивного интерфейса из утилитарных классов \cite{tailwind}. Такой подход удобен для страниц с большим количеством параметров генерации: размеры, отступы, состояния кнопок и сетки можно задавать непосредственно рядом с разметкой компонента.

Axios используется как HTTP-клиент для обращений frontend к API gateway \cite{axios}. Через него выполняются запросы регистрации и входа, получение списка фракталов, запуск генерации и загрузка изображений. Для прогресса рендера применяется поток событий \texttt{text/event-stream}.

\subsubsection{Система сборки Gradle}

Gradle используется для сборки Java-сервисов gateway и generator. Он обеспечивает управление зависимостями, генерацию Java-кода из protobuf-контрактов, запуск тестов и сборку исполняемых приложений. В gateway дополнительно подключён OpenAPI Generator: из файла \texttt{openapi.yaml} формируются Java-модели и API-интерфейсы, которые затем реализуются контроллерами сервиса \cite{gradle}.

\subsubsection{База данных PostgreSQL}

\begin{figure}[H]
  \centering
  
\includegraphics[width=0.75\textwidth]{postgresql.png}
  \caption{PostgreSQL}
\end{figure}

Для хранения структурированных данных используется система управления базами данных PostgreSQL \cite{postgresql}. Она выбрана из-за поддержки транзакций, ограничений целостности, индексов и зрелой интеграции с backend-стеком. Конкретная схема данных generator и IAM описана выше в подразделах о моделях данных, поэтому в данном месте фиксируется только обоснование выбора СУБД.

\subsubsection{MinIO и S3-совместимое хранилище}

Для хранения изображений используется MinIO --- S3-совместимое объектное хранилище \cite{minio}. Оно позволяет отделить бинарные PNG-файлы от реляционных данных и оставить PostgreSQL для метаданных, связей и поиска. Благодаря совместимости с S3 при необходимости локальное хранилище можно заменить внешним объектным хранилищем без изменения пользовательских сценариев.

\subsubsection{gRPC}

gRPC применяется для внутреннего взаимодействия gateway с generator и IAM \cite{grpc}. По сравнению с передачей JSON между backend-сервисами protobuf-контракты дают строгую схему сообщений, компактное представление данных и удобную генерацию клиентского и серверного кода. Подробный сценарий потоковой передачи прогресса рассмотрен в подразделе о внутренних gRPC-контрактах.

\subsubsection{JWT}

JSON Web Token используется как формат access token \cite{jwt_rfc}. Токен содержит идентификатор пользователя и имя, подписывается секретным ключом и проверяется без обращения к базе данных при каждом запросе. Благодаря этому gateway остаётся stateless, а механизм передачи пользовательского контекста между сервисами не требует общей HTTP-сессии.

\subsubsection{Контейнеризация Docker и Docker Compose}

Для развёртывания использовалась контейнеризация Docker. Каждый микросервис упакован в отдельный образ, а общий сценарий запуска описан в Docker Compose \cite{docker_compose}. Compose выбран как достаточный для курсового проекта способ поднять приложение, базу данных и MinIO одной командой в воспроизводимом локальном или серверном окружении.

\subsubsection{Среда разработки Visual Studio Code}

\begin{figure}[H]
  \includegraphics[width=\textwidth]{vscode.png}
  \caption{Экран редактирования кода Visual Studio Code}
\end{figure}

Visual Studio Code (VS Code) — это лёгкая, но мощная интегрированная среда разработки, разработанная компанией Microsoft. С момента своего выхода в 2015 году она завоевала огромную популярность благодаря простоте, быстроте и гибкости. VS Code является кроссплатформенным инструментом, доступным на Windows, macOS и Linux, что делает его универсальным решением для разработчиков по всему миру. \cite{vscode}

VS Code не перегружена лишними функциями, но в то же время предоставляет всё необходимое для разработки: интеллектуальную подсветку синтаксиса, автодополнение, отладку, поддержку Git, систему контроля версий и встроенный терминал. Это позволяет разрабатывать и тестировать код быстро и эффективно, без лишних задержек.

Основные причины выбора VS Code включают:
\begin{enumerate}
  \item
        Лёгкость и быстродействие. VS Code быстро запускается и работает с минимальными задержками, что делает её идеальной для разработки фронтенд-части и быстрого прототипирования.
  \item
        Расширяемость. Благодаря множеству доступных расширений, VS Code легко адаптируется под любые задачи.
  \item
        Интеграция с Git и другими системами контроля версий. VS Code имеет встроенную поддержку Git, что упрощает работу с репозиториями и позволяет эффективно управлять версиями проекта.
\end{enumerate}

Разработка всех компонентов проекта выполнялась исключительно в Visual Studio Code, что обеспечило единообразие среды и повторяемость процесса. Дополнительно использовались расширения для Java, Docker и работы с \LaTeX, что позволило вести и разработку, и оформление документации в одном инструменте. Даже этот документ написан в VS Code с использованием \LaTeX.

\section{Тестирование программного обеспечения}

\subsection{Стратегия тестирования}

Тестирование системы проводилось на нескольких уровнях: модульном, интеграционном и системном. На модульном уровне проверялась корректность отдельных компонентов генерации и обработки параметров. Интеграционные тесты подтверждали корректность взаимодействия микросервисов, а системные тесты моделировали реальную работу пользователя. Подход соответствует классическим рекомендациям по организации тестирования программного обеспечения \cite{testing_myers}.

Для проверки критических сценариев использовались контрольные наборы параметров рендера и заранее подготовленные эталонные результаты. Это позволило выявлять регрессии при изменениях алгоритмов и конфигурации.

\subsection{Сценарии тестирования}

Ниже приведены основные сценарии, подтверждающие работоспособность системы:
\begin{enumerate}
  \item Регистрация и вход пользователя: создание записи в IAM, получение access token и refresh token.
  \item Проверка JWT: передача access token из frontend в gateway и дальнейшая передача токена во внутренний gRPC-вызов.
  \item Создание задачи генерации: отправка параметров фрактала через HTTP API gateway и получение идентификатора задачи.
  \item Подписка на прогресс: получение событий \texttt{text/event-stream}, сформированных gateway на основе gRPC stream generator.
  \item Сохранение результата: запись параметров фрактала в PostgreSQL и PNG-изображения в MinIO.
  \item Просмотр истории: получение списка фракталов с пагинацией, сортировкой и фильтрацией по пользователю.
  \item Получение изображения: запрос изображения через gateway и чтение объекта из MinIO сервисом generator.
  \item Проверка ограничений параметров: отказ generator при превышении максимальных размеров или количества итераций.
\end{enumerate}

Результаты тестирования подтвердили корректность работы ключевых функций и соответствие заявленным требованиям.

Отдельно фиксировались временные показатели рендера и устойчивость к отказам отдельных сервисов, что подтвердило работоспособность системы в условиях повышенной нагрузки.

\newpage

\addcontentsline{toc}{section}{Заключение}
\section*{\centering{\uppercase{Заключение}}}

В результате выполнения курсового проекта получена проектная и программная основа веб-системы для генерации изображений фрактального пламени. Система закрывает основной недостаток рассмотренных аналогов: вместо локального desktop-инструмента пользователь получает веб-интерфейс с запуском генерации, просмотром прогресса, сохранением результата, общей галереей последних работ и личной историей. Примеры итоговой галереи фракталов вынесены в приложение~А.

Главным инженерным результатом стала микросервисная архитектура из четырёх компонентов: frontend, API gateway, IAM и generator. Frontend работает с gateway по HTTP, а gateway взаимодействует с IAM и generator по gRPC. Такое разделение позволило отделить пользовательский интерфейс от внутренней логики аутентификации и вычислений. Внешний API зафиксирован OpenAPI-спецификацией, а внутренний контракт generator описан protobuf-схемой, что уменьшает риск расхождения между клиентом, gateway и backend-сервисами.

Отдельным результатом стала реализация сервиса generator. Он принимает параметры фрактального пламени, проверяет ограничения, сохраняет модель фрактала, ставит задачу в очередь, выполняет многопоточный рендер, применяет гамма-коррекцию и сохраняет PNG-изображение в MinIO. Метаданные результата сохраняются в PostgreSQL, а связь с пользователем выполняется через идентификатор из JWT. Код генератора и сервиса генерации должен быть приведён в приложении~Б.

Для пользовательского опыта был решён сценарий длительной операции генерации. Generator хранит состояние задачи, scheduler периодически формирует обновления прогресса, gRPC stream передаёт их в gateway, а gateway преобразует поток во внешний SSE-канал для браузера. Благодаря этому пользователь видит не только факт запуска задачи, но и промежуточный предпросмотр изображения вместе с числовым прогрессом.

Также получены результаты по хранению и развёртыванию. PostgreSQL используется для структурированных данных IAM и generator, MinIO --- для изображений, Docker Compose --- для воспроизводимого запуска сервисов и инфраструктуры. Подготовленная схема развёртывания позволяет запускать один и тот же набор компонентов локально и на VPS, меняя только переменные окружения и настройки прокси.

Таким образом, цель курсового проекта достигнута: спроектирована и реализована основа распределённого приложения, поддерживающего регистрацию и вход пользователя, генерацию фракталов, передачу прогресса, сохранение истории, выдачу изображений и просмотр результатов через веб-интерфейс. Полученное решение остаётся расширяемым: в дальнейшем можно добавить новые вариационные функции, роли пользователей, очередь фоновых задач, более гибкий редактор параметров и дополнительные способы публикации готовых изображений.
